\chapter{Primary Fock Determinant: Rank Residue \(2\)}
\label{v4-arith:w9-r3:chapter}

\emph{Chapter~\ref{v4-arith:polyakov} separates the local conformal
charge from the global arithmetic determinant. The naive sum
\(\sum_{p}c^{(p)}\) diverges, and the \(\log p\)-weighted arithmetic
Hamiltonian is not a Virasoro zero mode. The determinant assertion is
the following: the primary Fock Euler product has
rank residue \(2\) at \(s=1\). It does not construct a global Virasoro
action on \(\mathcal{V}^{\mathrm{prim}}\), and it does not close
\Cref{v4-arith:rh:conj:primary-Virasoro}.}

\section{Local Virasoro and Arithmetic Hamiltonian}
\label{v4-arith:w9-r3:local}

\subsection{The per-prime chiral Heisenberg Fock}
\label{v4-arith:w9-r3:per-prime-fock}

\begin{definition}[per-prime Heisenberg and arithmetic Hamiltonian]
\label{v4-arith:w9-r3:def:per-prime}
\ClaimStatusDefinitional. At each finite place \(p\in\mathbb{P}\) the
\emph{per-prime chiral Heisenberg} is the complex rank-2 Heisenberg
vertex algebra \(\mathcal{H}^{(p)}\) with generators \(a^{(p)},b^{(p)}\), OPE
\(a^{(p)}(z)a^{(p)}(w)\sim (z-w)^{-2}\),
\(b^{(p)}(z)b^{(p)}(w)\sim (z-w)^{-2}\), mixed OPE vanishing, and Fock
vacuum \(|0\rangle_{p}\)
(Chapter~\ref{v4-arith:drinfeld-kapranov}; Kac 1998, \S 3.5).
The local conformal vector is
\begin{equation}
\label{v4-arith:w9-r3:eq:omega-p}
\omega_{p}^{\mathrm{conf}}\;:=\;\tfrac{1}{2}\bigl(a_{-1}^{(p)}a_{-1}^{(p)}+b_{-1}^{(p)}b_{-1}^{(p)}\bigr)|0\rangle_{p}.
\end{equation}
The state-field correspondence produces conformal modes
\begin{equation}
\label{v4-arith:w9-r3:eq:L-m-p}
L_{m}^{\mathrm{conf},(p)}\;=\;\tfrac{1}{2}\sum_{k\in\mathbb{Z}}{:}a_{m-k}^{(p)}a_{k}^{(p)}{:}\;+\;\tfrac{1}{2}\sum_{k\in\mathbb{Z}}{:}b_{m-k}^{(p)}b_{k}^{(p)}{:}.
\end{equation}
The arithmetic Hamiltonian is
\begin{equation}
\label{v4-arith:w9-r3:eq:N-p}
N_{p}\;:=\;(\log p)L_{0}^{\mathrm{conf},(p)}.
\end{equation}
\end{definition}

\begin{remark}[convention: \(\log p\) is archimedean]
\label{v4-arith:w9-r3:rem:log-p-archimedean}
The factor \(\log p\) is archimedean-real scaling data, not an
internal \(\mathbb{Q}_{p}\)-coefficient
(\Cref{v4-arith:conv:log-p-archimedean}). Observation
\Cref{v4-arith:obs:rescaling-breaks-virasoro} shows that rescaling
\(T\mapsto\lambda T\) breaks the Virasoro OPE unless \(\lambda=1\).
Thus \(N_{p}\) is a Hamiltonian used in the Mellin character, not a
Virasoro generator.
\end{remark}

\subsection{Per-prime central charge}
\label{v4-arith:w9-r3:per-prime-central}

\begin{proposition}[per-prime central charge is \(2\)]
\label{v4-arith:w9-r3:prop:per-prime-c-2}
\ClaimStatusProvedHere. For every prime \(p\) and every pair \(m,n\)
with \(m+n=0\) the Virasoro commutator on \(\mathcal{H}^{(p)}\) satisfies
\begin{equation}
\label{v4-arith:w9-r3:eq:per-prime-bracket}
[L_{m}^{\mathrm{conf},(p)},L_{-m}^{\mathrm{conf},(p)}]\;=\;2m\,L_{0}^{\mathrm{conf},(p)}\;+\;\frac{c^{(p)}}{12}\,m(m^{2}-1)\cdot\mathrm{id},\qquad c^{(p)}\;=\;2.
\end{equation}
\end{proposition}

\begin{proof}
Two independent Sugawara-type Heisenbergs tensor multiplicatively at
the level of Virasoro central charge: the single chiral \(u(1)\)
Heisenberg has \(c=1\) by Kac--Raina \S 3.5 / Frenkel--Ben-Zvi
Chapter~2. The rank-2 Heisenberg is the product, so
\(c^{(p)}=1+1=2\). The arithmetic Hamiltonian \(N_{p}\) in
\eqref{v4-arith:w9-r3:eq:N-p} is external to this Virasoro
calculation.
\end{proof}

\begin{remark}[independence of \(p\)]
\label{v4-arith:w9-r3:rem:p-independence}
\Cref{v4-arith:thm:per-prime-c} records \(c^{(p)}=2\) as a Polyakov
observation. Proposition~\ref{v4-arith:w9-r3:prop:per-prime-c-2}
derives it from the Sugawara tensor rule. The \(\log p\)-weighted
Hamiltonian enters only the determinant.
\end{remark}

\section{Divergent Sum and Per-State Finiteness}
\label{v4-arith:w9-r3:per-state}

\subsection{The naive global sum diverges}
\label{v4-arith:w9-r3:full-divergent}

\begin{observation}[full global conformal central sum is \(+\infty\)]
\label{v4-arith:w9-r3:obs:full-divergent}
\ClaimStatusProvedHere. The naive \emph{full} sum of local conformal
central charges
\begin{equation}
\label{v4-arith:w9-r3:eq:c-full-divergent}
c_{\mathrm{full}}\;:=\;\sum_{p\in\mathbb{P}}c^{(p)}\;=\;\sum_{p\in\mathbb{P}}2\;=\;+\infty
\end{equation}
diverges. Equivalently, the formal stress tensor
\(\sum_{p}\omega_{p}^{\mathrm{conf}}\) is not a vector of the
restricted tensor product, and no Virasoro central extension with a
finite scalar acts on \(\mathcal{V}^{\mathrm{prim}}\)
(Chapter~\ref{v4-arith:polyakov}, \Cref{v4-arith:cor:no-virasoro}).
\end{observation}

\subsection{The primary core}
\label{v4-arith:w9-r3:primary-core}

\begin{definition}[primary Heisenberg core]
\label{v4-arith:w9-r3:def:primary-core}
\ClaimStatusDefinitional. The \emph{primary Heisenberg core} of
\(\mathcal{V}^{\mathrm{prim}}\) is
\begin{equation}
\label{v4-arith:w9-r3:eq:primary-core}
\mathcal{V}^{\mathrm{prim}}_{\mathrm{Heis},\mathrm{prim}}\;:=\;\mathrm{span}_{\mathbb{C}}\Bigl\{\;v=\bigotimes_{p}v_{p}\;\Big|\;v_{p}=|0\rangle_{p}\text{ for all but finitely many }p\;\Bigr\},
\end{equation}
with \(|0\rangle_{p}\in\mathcal{H}^{(p)}\) the per-prime Fock vacuum
(\Cref{v4-arith:w5-b:def:fock}). Each \(v\in\mathcal{V}^{\mathrm{prim}}_{\mathrm{Heis},\mathrm{prim}}\)
has a finite \emph{support set} \(\mathrm{supp}(v)\subset\mathbb{P}\),
\(|\mathrm{supp}(v)|<\infty\).
\end{definition}

\begin{remark}[primary = finite-support]
\label{v4-arith:w9-r3:rem:primary-is-fin-supp}
The primary core is the algebraic restricted tensor product, before
any Hilbert completion: the dense subspace on which every
per-prime mode acts by zero at all but finitely many places. It is a
\(\mathbb{C}\)-vector space (not a \(\mathbb{Q}_{p}\)-module) carrying the
\(\mathbb{N}^{\times}\)-action by Adams operations
\(\psi_{n}\) (Adams-Frobenius; \Cref{v4-arith:def:lambda-ring} and
\Cref{v4-arith:thm:lambda-ring-descent}).
\end{remark}

\subsection{Finite Positive Modes and Divergent Negative Modes}
\label{v4-arith:w9-r3:per-state-finite}

\begin{proposition}[finite positive modes]
\label{v4-arith:w9-r3:prop:per-state-finite}
\ClaimStatusProvedHere. For every \(v\in\mathcal{V}^{\mathrm{prim}}_{\mathrm{Heis},\mathrm{prim}}\)
and every \(m\geq -1\) the conformal mode action
\begin{equation}
\label{v4-arith:w9-r3:eq:finite-acting}
L_{m}^{\geq -1}v\;:=\;\sum_{p\in S_{v}(m)}L_{m}^{\mathrm{conf},(p)}v,\qquad S_{v}(m)\;\subset\;\mathbb{P},\quad |S_{v}(m)|<\infty,
\end{equation}
is a finite sum. For \(m\leq -2\), the sum
\(\sum_{p}L_{m}^{\mathrm{conf},(p)}\) is not defined on the vacuum.
Consequently the global Virasoro commutator is not an operator on the
restricted tensor product. Only the local commutators exist, and their
central scalar has divergent prime sum.
\end{proposition}

\begin{proof}
Fix \(v=\bigotimes_{p}v_{p}\) with \(\mathrm{supp}(v)\subset\mathbb{P}\)
finite. For \(p\notin\mathrm{supp}(v)\), \(v_{p}=|0\rangle_{p}\) and
\(L_{m}^{\mathrm{conf},(p)}|0\rangle_{p}=0\) for every \(m\geq-1\)
(Kac--Raina \S 3.5; the standard creation-annihilation normal
ordering places any positive-mode on the vacuum to zero, and the
\(L_{-1},L_{0}\) annihilation on vacuum is the standard vertex-algebra axiom).
Thus only finitely many summands act non-trivially for \(m\geq-1\).
For \(m\leq-2\), each vacuum factor contributes the non-zero descendant
\(L_{m}^{\mathrm{conf},(p)}|0\rangle_{p}\). The sum over primes is then
an infinite sum of mutually orthogonal one-prime descendants, not an
element of the algebraic restricted tensor product. The local identity
\[
[L_{m}^{\mathrm{conf},(p)},L_{-m}^{\mathrm{conf},(p)}]|0\rangle_{p}
=\frac{c^{(p)}}{12}m(m^{2}-1)|0\rangle_{p}
\]
therefore produces the divergent scalar
\(\sum_{p}c^{(p)}\) when summed over all primes. This scalar is not an
operator before regularisation.
\end{proof}

\begin{remark}[finite action, divergent scalar]
\label{v4-arith:w9-r3:rem:central-term-decomposition}
\Cref{v4-arith:w9-r3:prop:per-state-finite} proves only finite action
of the conformal modes \(m\geq -1\). The formal central term
\begin{equation}
\label{v4-arith:w9-r3:eq:central-term-sum}
\frac{1}{12}m(m^{2}-1)\cdot\Bigl(\sum_{p}c^{(p)}\Bigr)\cdot\mathrm{id}
\end{equation}
contains \(\sum_{p}c^{(p)}=\infty\). The residue calculation below
assigns a scalar to the determinant; it does not turn the formal
commutator into a Virasoro action.
\end{remark}

\section{Dirichlet-Series Regularisation}
\label{v4-arith:w9-r3:dirichlet}

\subsection{The naive prime-zeta attempt fails}
\label{v4-arith:w9-r3:naive-fails}

\begin{observation}[naive prime-zeta regularisation is not well-posed]
\label{v4-arith:w9-r3:obs:prime-zeta-fails}
\ClaimStatusProvedHere. Define the formal series
\begin{equation}
\label{v4-arith:w9-r3:eq:c-s-prime-zeta}
c_{\mathbb{P}}(s)\;:=\;\sum_{p\in\mathbb{P}}c^{(p)}p^{-(s-1)}\;=\;2\sum_{p}p^{-(s-1)}\;=\;2\zeta_{\mathbb{P}}(s-1).
\end{equation}
The prime zeta \(\zeta_{\mathbb{P}}\) has natural boundary on
\(\mathrm{Re}(w)=0\) by the Landau--Walfisz / Fr\"oberg theorem
(\Cref{v4-arith:thm:prime-zeta-natural-boundary}). Thus \(c_{\mathbb{P}}(s)\)
at \(s=1\) asks for \(2\zeta_{\mathbb{P}}(0)\), which is undefined
(\Cref{v4-arith:cor:zeta-P-0-undefined}); at \(s\to 1^{+}\),
\(c_{\mathbb{P}}(s)\to+\infty\) logarithmically.
\end{observation}

\begin{remark}[prime-zeta sum vs.\ Riemann-zeta residue]
\label{v4-arith:w9-r3:rem:prime-zeta-error}
The Dirichlet series \(c_{\mathbb{P}}(s)\) uses the prime-index set
\(\mathbb{P}\) as its summation domain. The determinant-rank scalar is
instead read from the Euler product for the Riemann zeta
\(\zeta(s)\), which is meromorphic on \(\mathbb{C}\). The prime zeta
\(\zeta_{\mathbb{P}}\) has a natural boundary at
\(\mathrm{Re}(w)=0\), and therefore cannot define the value at
\(w=0\). The finite-mode statement
\Cref{v4-arith:w9-r3:prop:per-state-finite} gives only an algebraic
domain for the modes \(m\geq -1\); the residue at \(s=1\) belongs to
the determinant, not to a Virasoro commutator.
\end{remark}

\subsection{Arithmetic black-hole motivation: primon gas as the first determinant}
\label{v4-arith:w9-r3:black-hole-primon-motivation}

The partition function \(Z(s)=\prod_{p}\prod_{n\geq 1}(1-p^{-sn})^{-2}\)
defined in \S\ref{v4-arith:w9-r3:partition-reg} is the arithmetic-Chern--Simons
thermal determinant over prime knots before the zero-side trace formula is
available. The following heuristic identifies its origin in the primon-gas
thermodynamics of Julia; no black-hole existence theorem enters the proof of
Lemma~\ref{v4-arith:w9-r3:lem:Z-is-zeta-prod}, and no positivity statement
is inferred from the calculation.

In ordinary three-dimensional gravity written as Chern--Simons
theory, a black-hole sector is not a propagating local graviton. It is
a flat connection with non-trivial holonomy around the cycles of a
solid torus. Its one-loop determinant is built from the closed
geodesics of the filling. The arithmetic replacement used throughout
Chapters~\ref{v4-arith:3d-acs-E1bar}--\ref{v4-arith:w9-r1:wilson-frobenius}
is:
\begin{equation}
\label{v4-arith:w9-r3:eq:black-hole-dictionary}
\text{closed geodesic}\quad\rightsquigarrow\quad
\text{prime knot }K_{p},\qquad
\ell(K_{p})=\log p,\qquad
\mathrm{hol}(K_{p})=\mathrm{Frob}_{p}.
\end{equation}
At the level relevant here, the arithmetic Chern--Simons thermal
determinant is computed around a holonomy saddle, with the primitive
closed orbits indexed by primes.

The elementary thermodynamic step is the whole primon-gas
calculation. A single bosonic occupation mode over the orbit \(K_{p}\)
has energy
\begin{equation}
\label{v4-arith:w9-r3:eq:primon-energy}
E_{p}=\ell(K_{p})=\log p .
\end{equation}
At inverse temperature \(s\), its local partition function is the
geometric series
\begin{equation}
\label{v4-arith:w9-r3:eq:single-primon-factor}
\sum_{m\geq0}e^{-smE_{p}}
\;=\;\sum_{m\geq0}p^{-sm}
\;=\;(1-p^{-s})^{-1}.
\end{equation}
Multiplying the independent local factors over all prime knots gives
Julia's primon gas:
\begin{equation}
\label{v4-arith:w9-r3:eq:primon-zeta}
Z_{\mathrm{primon}}(s)
\;=\;\prod_{p}(1-p^{-s})^{-1}
\;=\;\zeta(s),\qquad \mathrm{Re}(s)>1.
\end{equation}
For a rank-\(r\) free Heisenberg sector the same calculation gives
\(\zeta(s)^{r}\). Here \(r=2\), because the local
arithmetic Heisenberg is generated by two commuting \(u(1)\) fields.

The oscillator refinement adds the usual Fock levels. Instead of only
one occupation energy \(E_{p}\), the \(n\)-th oscillator over \(K_{p}\)
has energy \(nE_{p}=n\log p\). Therefore the rank-\(2\) local factor is
\begin{equation}
\label{v4-arith:w9-r3:eq:oscillator-primon-factor}
\prod_{n\geq1}(1-e^{-s n\log p})^{-2}
\;=\;\prod_{n\geq1}(1-p^{-sn})^{-2}.
\end{equation}
Multiplication over all primes gives exactly the global Fock
partition function defined below:
\begin{equation}
\label{v4-arith:w9-r3:eq:black-hole-to-fock}
\prod_{p}\prod_{n\geq1}(1-p^{-sn})^{-2}.
\end{equation}
The primon gas is the \(n=1\) factor in the oscillator product: setting
\(n=1\) in \eqref{v4-arith:w9-r3:eq:oscillator-primon-factor} and taking
\(r=1\) recovers \((1-p^{-s})^{-1}\). The full rank-\(2\) Fock product
\eqref{v4-arith:w9-r3:eq:black-hole-to-fock} is the arithmetic
Chern--Simons thermal determinant used for the determinant
rank-residue extraction in \S\ref{v4-arith:w9-r3:partition-reg}.

The limitation is as important as the calculation. Equations
\eqref{v4-arith:w9-r3:eq:primon-zeta}--\eqref{v4-arith:w9-r3:eq:black-hole-to-fock}
recover the prime-side Euler product in its convergence half-plane.
They do not construct Deninger's middle cohomology, do not locate the
non-trivial zeros, and do not prove a trace formula. Those require the
spectral-flow and positivity inputs isolated in
Chapters~\ref{v4-arith:w8-a:deninger-chiral-homology} and
\ref{v4-arith:tfe-gutzwiller}.

\subsection{The partition function regularisation}
\label{v4-arith:w9-r3:partition-reg}

\begin{definition}[per-prime Fock partition]
\label{v4-arith:w9-r3:def:per-prime-partition}
\ClaimStatusDefinitional. For each prime \(p\) and complex parameter
\(s\) with \(\mathrm{Re}(s)>0\) the \emph{per-prime Fock partition} is
\begin{equation}
\label{v4-arith:w9-r3:eq:per-prime-Z}
Z^{(p)}(s)\;:=\;\mathrm{Tr}_{\mathcal{H}^{(p)}}\bigl(e^{-sN_{p}}\bigr)
\;=\;\mathrm{Tr}_{\mathcal{H}^{(p)}}\bigl(p^{-sL_{0}^{\mathrm{conf},(p)}}\bigr)
\;=\;\prod_{n\geq 1}(1-p^{-sn})^{-2},
\end{equation}
the final equality holding formally because \(L_{0}^{\mathrm{conf},(p)}\)
on the rank-\(2\) Heisenberg Fock has generating function \(\eta\)-type with
two commuting copies (Chapter~\ref{v4-arith:polyakov}, \S
\ref{v4-arith:pol:c-as-residue}; Kac--Raina \S 3.5 for the rank-1
generating series).
\end{definition}

\begin{definition}[global Fock partition: Euler product]
\label{v4-arith:w9-r3:def:global-partition}
\ClaimStatusDefinitional. The \emph{global Fock partition} is the
formal Euler product
\begin{equation}
\label{v4-arith:w9-r3:eq:global-Z}
Z(s)\;:=\;\prod_{p\in\mathbb{P}}Z^{(p)}(s)\;=\;\prod_{p}\prod_{n\geq 1}(1-p^{-sn})^{-2},
\end{equation}
convergent for \(\mathrm{Re}(s)>1\) by convergence of the log-sum
\(\sum_{p,n}p^{-sn}\) (standard Euler-product estimate, Serre 1973
Ch.~VI).
\end{definition}

\begin{lemma}[partition function is squared zeta product]
\label{v4-arith:w9-r3:lem:Z-is-zeta-prod}
\ClaimStatusProvedHere. In the region \(\mathrm{Re}(s)>1\),
\begin{equation}
\label{v4-arith:w9-r3:eq:Z-is-zeta-prod}
Z(s)\;=\;\prod_{n\geq 1}\zeta(ns)^{2}.
\end{equation}
\end{lemma}

\begin{proof}
Interchange the two products in \eqref{v4-arith:w9-r3:eq:global-Z}:
\[
Z(s)\;=\;\prod_{n\geq 1}\prod_{p}(1-p^{-sn})^{-2}\;=\;\prod_{n\geq 1}\Bigl(\prod_{p}(1-p^{-sn})^{-1}\Bigr)^{2}\;=\;\prod_{n\geq 1}\zeta(ns)^{2}.
\]
Each inner factor is the standard Euler product of \(\zeta(ns)\) for
\(\mathrm{Re}(ns)>1\), valid since \(\mathrm{Re}(s)>1\) implies
\(\mathrm{Re}(ns)>1\) for every \(n\geq 1\). The interchange of infinite
products is justified by absolute convergence of the double sum
\(\sum_{p,n}p^{-\mathrm{Re}(s)n}\) on compact subsets of
\(\{\mathrm{Re}(s)>1\}\).
\end{proof}

\begin{remark}[logarithmic derivative of \(Z\) and the rank pole]
\label{v4-arith:w9-r3:rem:log-Z-central}
From \(\log Z(s)=2\sum_{n\geq 1}\log\zeta(ns)\), the logarithmic derivative is
\begin{equation}
\label{v4-arith:w9-r3:eq:log-Z-prime}
-\frac{d}{ds}\log Z(s)\;=\;2\sum_{n\geq 1}\frac{n\,\zeta'(ns)}{\zeta(ns)}\cdot(-1)\;=\;2\sum_{n\geq 1}n\,\frac{-\zeta'(ns)}{\zeta(ns)},
\end{equation}
convergent for \(\mathrm{Re}(s)>1\). The leading pole of \(-\frac{d}{ds}\log Z\) at
\(s=1\) arises from the \(n=1\) term via \(\frac{-\zeta'(s)}{\zeta(s)}\sim (s-1)^{-1}\)
near \(s=1\), with the higher-\(n\) terms holomorphic at \(s=1\).
In the corrected Polyakov convention
(\S\ref{v4-arith:pol:c-as-residue}), this residue is the rank scalar
of the determinant, not a global central charge.
\end{remark}

\subsection{The Dirichlet-series rank residue}
\label{v4-arith:w9-r3:dirichlet-reg}

\begin{definition}[zeta rank residue]
\label{v4-arith:w9-r3:def:c-reg}
\ClaimStatusDefinitional. The \emph{zeta rank residue} is
\begin{equation}
\label{v4-arith:w9-r3:eq:c-reg-def}
r_{\zeta}\;:=\;\mathrm{Res}_{s=1}\Bigl(-\frac{d}{ds}\log Z(s)\Bigr).
\end{equation}
Equivalently, if \(\log Z(s)=-r\log(s-1)+\mathrm{holomorphic}\) near \(s=1\),
then \(r_{\zeta}:=r\), where \(r\) is the order of the logarithmic
singularity.
\end{definition}

\begin{remark}[rationale for the residue definition]
\label{v4-arith:w9-r3:rem:residue-rationale}
Per Polyakov (Chapter~\ref{v4-arith:polyakov},
\Cref{v4-arith:thm:regularised-c}), the honest
regularisation is \emph{not} a sum-over-primes but the residue of
the global partition-function pole. The definition
\eqref{v4-arith:w9-r3:eq:c-reg-def} places this residue at the
\(s=1\) pole of \(\zeta\), which is the dominant pole of \(Z\) since
\(Z(s)=\zeta(s)^{2}\cdot\prod_{n\geq 2}\zeta(ns)^{2}\) and the \(n\geq 2\)
factors are holomorphic at \(s=1\) (with \(\zeta(2n)\in\mathbb{Q}\pi^{2n}\)
for \(n\geq 1\)).
\end{remark}

\section{Local Expansion at \(s=1\)}
\label{v4-arith:w9-r3:continuation}

\subsection{Decomposition of \(\log Z\) near \(s=1\)}
\label{v4-arith:w9-r3:log-Z-decomp}

\begin{theorem}[local Laurent expansion of \(\log Z\) at \(s=1\)]
\label{v4-arith:w9-r3:thm:log-Z-laurent}
\ClaimStatusProvedHere. In a punctured neighbourhood of \(s=1\) on a
fixed logarithm branch, the function
\(\log Z(s)=2\sum_{n\geq 1}\log\zeta(ns)\) has expansion
\begin{equation}
\label{v4-arith:w9-r3:eq:log-Z-laurent}
\log Z(s)\;=\;-2\log(s-1)\;+\;2\sum_{n\geq 2}\log\zeta(n)
\;+\;2\left(\gamma_{E}+\sum_{n\geq2}n\,{\zeta'(n)\over\zeta(n)}\right)(s-1)
\;+\;O((s-1)^{2}),
\end{equation}
where \(\gamma_{E}=\lim_{N\to\infty}(\sum_{k=1}^{N}1/k-\log N)\) is
the Euler--Mascheroni constant.
\end{theorem}

\begin{proof}
Write \(\log Z(s)=2\log\zeta(s)+2\sum_{n\geq 2}\log\zeta(ns)\). For
\(n\geq 2\) and \(s\) near \(1\), \(ns\geq 2\), so \(\zeta(ns)\) is
holomorphic and nonzero at \(s=1\) (since
\(\zeta(2),\zeta(3),\ldots\) are positive real and greater than \(1\)).
The tail bound \(\sum_{n\geq 2}|\log\zeta(ns)|\leq C\sum_{n\geq 2}2^{-n\mathrm{Re}(s)}\)
is uniform on \(\mathrm{Re}(s)\geq 1/2+\epsilon\), so \(\sum_{n\geq 2}\log\zeta(ns)\)
is holomorphic at \(s=1\) and has first derivative
\(\sum_{n\geq2}n\,\zeta'(n)/\zeta(n)\) there.

For the \(n=1\) term: the Laurent expansion \(\zeta(s)=(s-1)^{-1}+\gamma_{E}+O(s-1)\)
(Edwards 1974, \S 1.6) gives
\[
\log\zeta(s)=\log\bigl((s-1)^{-1}(1+\gamma_{E}(s-1)+O((s-1)^{2}))\bigr)
=-\log(s-1)+\gamma_{E}(s-1)+O((s-1)^{2}),
\]
using \(\log(1+x)=x+O(x^{2})\). Therefore
\[
\log Z(s)
=-2\log(s-1)+2\sum_{n\geq 2}\log\zeta(n)
+2\left(\gamma_{E}+\sum_{n\geq2}n\,{\zeta'(n)\over\zeta(n)}\right)(s-1)
+O((s-1)^{2}),
\]
which is \eqref{v4-arith:w9-r3:eq:log-Z-laurent}.
\end{proof}

\begin{remark}[why the rank residue is \(2\)]
\label{v4-arith:w9-r3:rem:residue-is-2}
The leading term \(-2\log(s-1)\) in \eqref{v4-arith:w9-r3:eq:log-Z-laurent}
has a \(-2\) coefficient because \(Z(s)=\zeta(s)^{2}\cdot\mathrm{hol}\)
near \(s=1\), and the double factor in \(\zeta\) comes from the
rank-\(2\) Heisenberg: per-prime two-commuting-Heisenberg implies
per-prime \(Z^{(p)}\) has a \emph{squared} \((1-p^{-s})^{-1}\) factor
leading at \(n=1\), and on the Euler-product reshuffle this makes
\(Z=\zeta^{2}\cdot\mathrm{hol}\). The power of \(\zeta\) at \(s=1\) is
the per-prime rank, here \(2\). This rank agrees with the local
Sugawara central charge prime by prime, but it is not a global central
charge.
\end{remark}

\subsection{Residue extraction}
\label{v4-arith:w9-r3:residue}

\begin{theorem}[residue at \(s=1\)]
\label{v4-arith:w9-r3:thm:residue-at-1}
\ClaimStatusProvedHere. With the decomposition
\eqref{v4-arith:w9-r3:eq:log-Z-laurent},
\begin{equation}
\label{v4-arith:w9-r3:eq:residue-extraction}
\mathrm{Res}_{s=1}\Bigl(-\frac{d}{ds}\log Z(s)\Bigr)\;=\;2,
\end{equation}
and the zeta rank residue of
\Cref{v4-arith:w9-r3:def:c-reg} evaluates to
\begin{equation}
\label{v4-arith:w9-r3:eq:c-reg-is-2}
r_{\zeta}\;=\;2.
\end{equation}
\end{theorem}

\begin{proof}
Differentiate \eqref{v4-arith:w9-r3:eq:log-Z-laurent}:
\[
-\frac{d}{ds}\log Z(s)=\frac{2}{s-1}
-2\left(\gamma_{E}+\sum_{n\geq2}n\,{\zeta'(n)\over\zeta(n)}\right)
+O(s-1)
\]
near \(s=1\). The residue at \(s=1\) is \(2\).
By Definition~\ref{v4-arith:w9-r3:def:c-reg},
\(r_{\zeta}=\mathrm{Res}_{s=1}(-\frac{d}{ds}\log Z)=2\).
\end{proof}

\begin{corollary}[\(r_{\zeta}=2\cdot\mathrm{Res}_{s=1}\zeta(s)\)]
\label{v4-arith:w9-r3:cor:c-reg-as-residue}
\ClaimStatusProvedHere. Since \(\mathrm{Res}_{s=1}\zeta(s)=1\),
\begin{equation}
\label{v4-arith:w9-r3:eq:c-reg-polyakov-form}
r_{\zeta}\;=\;2\cdot\mathrm{Res}_{s=1}\zeta(s)\;=\;2.
\end{equation}
\end{corollary}

\begin{remark}[Dirichlet-series form of regularisation]
\label{v4-arith:w9-r3:rem:dirichlet-form}
The original informal heuristic was
\(c(s)=\sum_{p}c^{(p)}p^{-(s-1)}=2\zeta_{\mathbb{P}}(s-1)\); we
replaced it by the honest global-partition-function residue. The
replacement is forced by
\Cref{v4-arith:thm:prime-zeta-natural-boundary}: the prime zeta
has natural boundary on \(\mathrm{Re}(w)=0\), so the informal
prescription has no analytic continuation to \(s=1\). The
partition-function prescription runs on \(\zeta(s)\) directly (which
is meromorphic on \(\mathbb{C}\)), with residue \(1\) at \(s=1\); the
total per-prime rank \(2\) and the partition-function-residue
identity \eqref{v4-arith:w9-r3:eq:c-reg-polyakov-form} together
deliver \(r_{\zeta}=2\). This is the determinant-rank residue of
\Cref{v4-arith:pol:c-as-residue}, not a sum over primes and not a
Virasoro central charge.
\end{remark}

\section{Polyakov Anomaly Connection}
\label{v4-arith:w9-r3:polyakov-connection}

\subsection{The Polyakov stress tensor and anomaly}
\label{v4-arith:w9-r3:polyakov-stress}

\begin{theorem}[Polyakov anomaly on the archimedean chiral Heisenberg]
\label{v4-arith:w9-r3:thm:polyakov-anomaly}
\ClaimStatusProvedElsewhere (Polyakov 1981; Belavin--Polyakov--Zamolodchikov
1984; Friedan--Shenker 1987; archimedean Heisenberg: Kac 1998
\S 3.5). On the archimedean Heisenberg \(\mathcal{H}^{(\infty)}\) with
rank-\(2\) free-field presentation, the Polyakov conformal anomaly
coefficient is
\begin{equation}
\label{v4-arith:w9-r3:eq:polyakov-anomaly}
c^{(\infty)}_{\mathrm{anomaly}}\;=\;2,
\end{equation}
extracted from the \(T(z)T(w)\) OPE central term with
\(T=\tfrac{1}{2}(a_{-1}^{(\infty)}a_{-1}^{(\infty)}+b_{-1}^{(\infty)}b_{-1}^{(\infty)})\),
independently from trace-anomaly \(\langle T^{\mu}_{\mu}\rangle=-cR/24\pi\)
and from Friedan--Shenker modular-invariance bootstrap.
\end{theorem}

\begin{remark}[archimedean vs.\ determinant]
\label{v4-arith:w9-r3:rem:archimedean-vs-full}
\(c^{(\infty)}_{\mathrm{anomaly}}=2\) is the anomaly of the
\emph{archimedean} Heisenberg alone. Each per-prime place also satisfies
\(c^{(p)}=2\), so the per-place charges are uniformly \(2\). The
finite-prime determinant residue also equals \(2\). Equality of these
two numbers is a rank comparison; it is not an addition of local
central charges.
\end{remark}

\subsection{The Polyakov--regularisation identity}
\label{v4-arith:w9-r3:polyakov-id}

\begin{theorem}[Polyakov-rank comparison]
\label{v4-arith:w9-r3:thm:polyakov-identity}
\ClaimStatusProvedHere. The zeta rank residue of
\Cref{v4-arith:w9-r3:def:c-reg} equals the Polyakov conformal anomaly
coefficient of the archimedean rank-\(2\) Heisenberg:
\begin{equation}
\label{v4-arith:w9-r3:eq:polyakov-id}
r_{\zeta}\;=\;c^{(\infty)}_{\mathrm{anomaly}}\;=\;2.
\end{equation}
Equivalently, the residue of the global Fock partition
\(Z(s)=\prod_{p}\prod_{n\geq 1}(1-p^{-sn})^{-2}\) at \(s=1\)
has the same value as the rank-\(2\) Polyakov-Friedan-Shenker anomaly
computed on the archimedean completion \(\mathcal{H}^{(\infty)}\).
\end{theorem}

\begin{proof}
Combine \Cref{v4-arith:w9-r3:cor:c-reg-as-residue} (\(r_{\zeta}=2\))
with \Cref{v4-arith:w9-r3:thm:polyakov-anomaly}
(\(c^{(\infty)}_{\mathrm{anomaly}}=2\)): both sides equal \(2\)
independently, and their common value is
\(2\cdot\mathrm{Res}_{s=1}\zeta(s)=2\).

The two sides arise from different constructions: the left-hand side
from a global partition-function Euler product over all finite primes,
regularised through the residue of \(\zeta\) at its pole; the
right-hand side from a Lagrangian stress-tensor computation on the
archimedean fibre, regularised through Polyakov's short-distance anomaly. The identity is the arithmetic-branch
expression of Polyakov's principle in the residue form
(\Cref{v4-arith:thm:regularised-c}): the finite-prime determinant and
the archimedean free-field anomaly have the same rank \(2\). No global
Virasoro action on \(\mathcal{V}^{\mathrm{prim}}\) is produced.
\end{proof}

\begin{remark}[global compatibility]
\label{v4-arith:w9-r3:rem:global-compatibility}
Product formula phrasing: for the completed zeta
\(\xi_{\mathbb{R}}(s)=\tfrac{1}{2}s(s-1)\pi^{-s/2}\Gamma(s/2)\zeta(s)\),
the factor \(\pi^{-s/2}\Gamma(s/2)\) is the archimedean local
\(\Gamma\)-factor and \(\zeta(s)\) is the Euler product over finite
primes. The residue of \(\zeta\) at \(s=1\) is \(1\) by the standard
prime-number-theorem residue; under the functional equation
\(\xi(s)=\xi(1-s)\), this residue is mirrored to the residue at
\(s=0\). On the conformal side, the Polyakov anomaly on the
archimedean fibre equals the rank (\(=2\)) of the free-field
Heisenberg; on the Euler-product side, the residue of \(Z\) at \(s=1\)
equals the per-prime rank (\(=2\)). The identification in
\Cref{v4-arith:w9-r3:thm:polyakov-identity} is rank \(=\) rank,
made precise through the \(\zeta\)-residue dictionary.
\end{remark}

\subsection{Derivation and Residue Sources}
\label{v4-arith:w9-r3:independence}

\begin{remark}[disjoint derivation and residue sources]
\label{v4-arith:w9-r3:rem:disjoint-sources}
The local chiral calculation uses Kac--Raina \emph{Bombay Lectures on Highest
Weight Representations of Infinite-Dimensional Lie Algebras}
(World Scientific 1987), \S 3.5 (Sugawara central charge for
bosonic Fock); Kac \emph{Vertex Algebras for Beginners} 2nd ed.
(AMS 1998), Chapter 3 (rank-\(n\) Heisenberg vertex algebra and Virasoro
central charge formula).
The global residue calculation uses Edwards \emph{Riemann's Zeta Function}
(Academic Press 1974), \S 1.6 (Laurent expansion of \(\zeta\) at
\(s=1\)) for the residue-extraction half; Serre \emph{Cours
d'Arithm\'etique} (PUF 1970) Ch.~VI (Euler product) for the
global partition rewriting. The two source sets (Kac/Kac--Raina for
the local central charge, Edwards/Serre for the global residue)
are disjoint; the computed scalar (\(2\)) matches on both sides.
\end{remark}

\section{Main Theorem: \(r_{\zeta}=2\)}
\label{v4-arith:w9-r3:main}

\subsection{Statement}
\label{v4-arith:w9-r3:main-statement}

\begin{theorem}[primary Fock determinant has rank residue \(2\)]
\label{v4-arith:w9-r3:thm:main}
\ClaimStatusProvedHere. Let \(\mathcal{V}^{\mathrm{prim}}_{\mathrm{Heis},\mathrm{prim}}\)
denote the primary Heisenberg core of
\Cref{v4-arith:w9-r3:def:primary-core}. Then:
\begin{enumerate}
\item[\textnormal{(a)}] \emph{Finite positive modes.} For every
\(v\in\mathcal{V}^{\mathrm{prim}}_{\mathrm{Heis},\mathrm{prim}}\) and
every \(m\geq -1\), the mode action
\(L_{m}^{\geq -1}v=\sum_{p\in S_{v}(m)}L_{m}^{\mathrm{conf},(p)}v\)
is a finite sum over \(|S_{v}(m)|<\infty\) primes. The modes
\(m\leq -2\) do not form global operators on the restricted vacuum
(\Cref{v4-arith:w9-r3:prop:per-state-finite}).
\item[\textnormal{(b)}] \emph{Divergent conformal scalar.} The
formal Virasoro central-term coefficient
\(\sum_{p}c^{(p)}\) is \(+\infty\) as a naive sum
(\Cref{v4-arith:w9-r3:obs:full-divergent}).
\item[\textnormal{(c)}] \emph{Zeta rank residue.}
The Dirichlet-series residue of the scalar via the global
Fock partition \(Z(s)=\prod_{p}Z^{(p)}(s)\) yields
\begin{equation}
\label{v4-arith:w9-r3:eq:main-result}
r_{\zeta}\;=\;\mathrm{Res}_{s=1}\Bigl(-\frac{d}{ds}\log Z(s)\Bigr)\;=\;2.
\end{equation}
\item[\textnormal{(d)}] \emph{Polyakov-rank comparison.} The rank
residue equals the Polyakov conformal anomaly coefficient of the
archimedean rank-\(2\) Heisenberg
\(c^{(\infty)}_{\mathrm{anomaly}}=2\)
(\Cref{v4-arith:w9-r3:thm:polyakov-identity}).
\end{enumerate}
\end{theorem}

\begin{proof}
(a) is \Cref{v4-arith:w9-r3:prop:per-state-finite}. (b) is
\Cref{v4-arith:w9-r3:obs:full-divergent}. (c) is
\Cref{v4-arith:w9-r3:thm:residue-at-1} with
\Cref{v4-arith:w9-r3:cor:c-reg-as-residue}. (d) is
\Cref{v4-arith:w9-r3:thm:polyakov-identity}.
\end{proof}

\subsection{Falsification and correction of six potential objections}
\label{v4-arith:w9-r3:obstruction-analysis}

\begin{remark}[prime-zeta sum does not regularise the determinant rank]
\label{v4-arith:w9-r3:obstruction-1}
\textbf{Objection.} ``Define \(c(s)=\sum_{p}c^{(p)}p^{-(s-1)}\), take
\(s\to 1\), get a finite number.''
\textbf{Answer.} \(c(s)=2\zeta_{\mathbb{P}}(s-1)\); \(\zeta_{\mathbb{P}}\)
has natural boundary on \(\mathrm{Re}(w)=0\)
(\Cref{v4-arith:thm:prime-zeta-natural-boundary}); \(c(s)\to+\infty\)
as \(s\to 1^{+}\). The correct regularisation is through the
determinant-residue prescription of
\Cref{v4-arith:w9-r3:def:c-reg}, which evaluates through
\(\zeta\) (meromorphic on \(\mathbb{C}\)) rather than \(\zeta_{\mathbb{P}}\)
(barrier at \(\mathrm{Re}(w)=0\)).
\end{remark}

\begin{remark}[\((\log p)\)-weighting is not a Virasoro rescaling]
\label{v4-arith:w9-r3:obstruction-2}
\textbf{Objection.} ``The \((\log p)\)-rescaling changes the central
charge \(c^{(p)}\mapsto(\log p)^{2}\cdot c^{(p)}\).''
\textbf{Answer.} The local stress tensor is not rescaled. The factor
\(\log p\) defines \(N_{p}=(\log p)L_{0}^{\mathrm{conf},(p)}\), the
Hamiltonian in the Mellin determinant. The local conformal charge
remains \(c^{(p)}=2\), and \(N_{p}\) is not a Virasoro zero mode.
\end{remark}

\begin{remark}[the universal scalar is not per-state-finite]
\label{v4-arith:w9-r3:obstruction-3}
\textbf{Objection.} ``\Cref{v4-arith:w9-r3:prop:per-state-finite}
establishes a global Virasoro commutator.''
\textbf{Answer.} It establishes finite action only for
\(m\geq -1\). For \(m\leq -2\) the global mode is not defined on the
vacuum, and the formal central scalar is \(\infty\). The residue in
\S\ref{v4-arith:w9-r3:dirichlet} targets the determinant scalar, not
a commutator on a dense domain
(\Cref{v4-arith:w9-r3:rem:central-term-decomposition}).
\end{remark}

\begin{remark}[the residue prescription is uniquely determined]
\label{v4-arith:w9-r3:obstruction-4}
\textbf{Objection.} ``The regularisation \(r_{\zeta}=2\) is
arbitrary: many Dirichlet-series prescriptions yield different
numbers.''
\textbf{Answer.} The prescription
\eqref{v4-arith:w9-r3:eq:c-reg-def} is fixed by two requirements:
(i) the regularisation runs through the global Fock partition
\(Z(s)\) whose pole structure is uniquely determined by the Euler
product of \(\zeta^{2}\) at \(s=1\); (ii) the regularisation must
match the Polyakov anomaly on the archimedean completion
(\Cref{v4-arith:w9-r3:thm:polyakov-identity}). Together these
single out the residue prescription; the prime-zeta prescription
fails (i): \(\zeta_{\mathbb{P}}\) has a natural boundary at \(\mathrm{Re}(w)=0\);
the sum-at-zero prescription fails (ii): a sum indexed by primes
is not the same as a residue of the Riemann zeta at \(s=1\), which
runs over all positive integers via the Euler product.
\end{remark}

\begin{remark}[the archimedean fibre enters through the completed-zeta \(\Gamma\)-factor]
\label{v4-arith:w9-r3:obstruction-5}
\textbf{Objection.} ``The partition \(Z(s)=\prod_{p}Z^{(p)}(s)\) omits
the archimedean fibre \(\mathcal{H}^{(\infty)}\).''
\textbf{Answer.} The archimedean contribution enters through the
\(\Gamma\)-factor in the completed zeta
\(\xi_{\mathbb{R}}(s)=\tfrac{1}{2}s(s-1)\pi^{-s/2}\Gamma(s/2)\zeta(s)\);
on the stress-tensor side it is built into the Polyakov-anomaly
identification. \Cref{v4-arith:w9-r3:thm:polyakov-identity}
states only a rank equality between the finite-prime residue and the
archimedean anomaly. It does not add two central charges.
\end{remark}

\begin{remark}[the order of the pole of \(\log Z\) at \(s=1\) is not a free parameter]
\label{v4-arith:w9-r3:obstruction-6}
\textbf{Objection.} ``The residue prescription is circular: the
residue is \(2\) because the rank is chosen to be \(2\).''
\textbf{Answer.} The order of the pole is \emph{not} a free parameter. It is
computed from the Euler product: \(Z^{(p)}(s)=\prod_{n\geq 1}(1-p^{-sn})^{-2}\)
has rank-\(2\) because the per-prime Heisenberg is rank-\(2\)
(two free bosons); this is a \emph{fixed input} from the
construction. The resulting \(Z(s)=\zeta(s)^{2}\cdot\mathrm{hol}\)
near \(s=1\) has a pole of order \(2\) in \(\zeta\), hence a leading
\(-2\log(s-1)\) term in \(\log Z\), hence
\(\mathrm{Res}_{s=1}(-\frac{d}{ds}\log Z)=+2\). The input is the
rank-\(2\) Heisenberg determinant. The output is the same rank read
from the pole of \(\zeta(s)^{2}\).
\end{remark}

\begin{remark}[Iwasawa completion does not change the local rank]
\label{v4-arith:w9-r3:obstruction-7}
\textbf{Objection.} ``The Vol~IV arithmetic-Positselski chapter
works with Iwasawa \(\Lambda_{p}\)-coefficients, where the local
rank of the Heisenberg factor might change.''
\textbf{Answer.} \(\Lambda_{p}=\mathbb{Z}_{p}[[T]]\) is the
Iwasawa-algebra completion of the per-prime Heisenberg acting on
the cyclotomic tower (Chapter~\ref{v4-arith:arith-positselski},
\Cref{v4-arith:thm:arithmetic-positselski}). In the
Gelfand--Kazhdan formal-geometry window the \(\Lambda_{p}\)-Heisenberg
is rank-\(2\) over \(\mathbb{Q}_{p}\) (two
free Iwasawa-generators; Fontaine 1991 \(B_{\mathrm{dR}}\) rank-\(2\)
Hodge-Tate class). Completion preserves this rank. Thus the local
Sugawara central charge remains \(2\), and the global trace class
still has no Virasoro central charge.
\end{remark}

\subsection{Relation to the conditional Virasoro-bracket question}
\label{v4-arith:w9-r3:bracket-relation}

\begin{remark}[what is proved and what is not]
\label{v4-arith:w9-r3:rem:proved-vs-not}
\Cref{v4-arith:w9-r3:thm:main} proves the determinant scalar
\(r_{\zeta}=2\). It does \emph{not} prove that a per-prime sum
\(\sum_{p}L_{m}^{\mathrm{conf},(p)}\) satisfies a Virasoro bracket
with this scalar. Two obstructions remain:
\begin{enumerate}
\item \emph{Domain.} The modes \(m\leq -2\) must first exist as
dense operators on a completion of
\(\mathcal{V}^{\mathrm{prim}}_{\mathrm{Heis},\mathrm{prim}}\) (the
Petersson completion, \Cref{v4-arith:thm:P3-resolution}). The
closure is not automatic: finite-support states do not form an
invariant core unless a commutator estimate is given.
\item \emph{Bracket identity.} After closure, the identity
\([L_{m},L_{n}]=(m-n)L_{m+n}+(c/12)(m^{3}-m)\delta_{m+n,0}\cdot\mathrm{id}\)
would still be a new theorem. The determinant scalar \(r_{\zeta}\)
does not supply it.
\end{enumerate}
Both obstructions are open for the primary-Virasoro reconciliation.
\end{remark}

\subsection{Domain and logical form}
\label{v4-arith:w9-r3:domain}

\begin{remark}[determinant domain]
\label{v4-arith:w9-r3:rem:domain}
The determinant-rank identity is \(r_{\zeta}=2\).
It does not prove either half of
\Cref{v4-arith:rh:conj:primary-Virasoro}. The open content remains
operator closure plus bracket identity on a dense domain
(\Cref{v4-arith:w9-r3:rem:proved-vs-not}). The operator closure
requires operator estimates on
\(\mathcal{V}^{\mathrm{prim}}_{\mathrm{Heis},\mathrm{prim}}\) with
respect to the Petersson inner product of
\Cref{v4-arith:thm:P3-resolution}.
\end{remark}

\section{Bibliography}
\label{v4-arith:w9-r3:bibliography}

\begin{thebibliography}{99}
\bibitem{Belavin-Polyakov-Zamolodchikov-1984}
A.~A.~Belavin, A.~M.~Polyakov, A.~B.~Zamolodchikov.
\emph{Infinite conformal symmetry in two-dimensional quantum field
theory.} Nuclear Physics B 241 (1984), 333--380.

\bibitem{Edwards-1974}
H.~M.~Edwards.
\emph{Riemann's Zeta Function.}
Pure and Applied Mathematics 58. Academic Press 1974.

\bibitem{Friedan-Shenker-1987}
D.~Friedan, S.~Shenker.
\emph{The analytic geometry of two-dimensional conformal field
theory.} Nuclear Physics B 281 (1987), 509--545.

\bibitem{Frenkel-Ben-Zvi-2004}
E.~Frenkel, D.~Ben-Zvi.
\emph{Vertex Algebras and Algebraic Curves.} 2nd ed.
Mathematical Surveys and Monographs 88. American Mathematical
Society 2004. Chapter 2 (Heisenberg vertex algebras), Chapter 5
(Sugawara construction).

\bibitem{Iwasawa-1972}
K.~Iwasawa.
\emph{Lectures on p-adic L-functions.}
Annals of Mathematics Studies 74. Princeton University Press 1972.

\bibitem{Kac-1998}
V.~G.~Kac.
\emph{Vertex Algebras for Beginners.} 2nd ed.
University Lecture Series 10. American Mathematical Society 1998.

\bibitem{Kac-Raina-1987}
V.~G.~Kac, A.~K.~Raina.
\emph{Bombay Lectures on Highest Weight Representations of
Infinite-Dimensional Lie Algebras.}
Advanced Series in Mathematical Physics 2. World Scientific 1987.

\bibitem{Landau-Walfisz-1920}
E.~Landau, A.~Walfisz.
\emph{\"Uber die Nichtfortsetzbarkeit einiger durch Dirichletsche
Reihen definierten Funktionen.}
Rendiconti del Circolo Matematico di Palermo 44 (1920), 82--86.

\bibitem{Polyakov-1981}
A.~M.~Polyakov.
\emph{Quantum geometry of bosonic strings.}
Physics Letters B 103 (1981), 207--210.

\bibitem{Serre-1970}
J.-P.~Serre.
\emph{Cours d'Arithm\'etique.}
Presses Universitaires de France 1970. Chapter VI (Dirichlet
series, Euler products, analytic continuation of \(\zeta\)).

\bibitem{Froberg-1968}
C.-E.~Fr\"oberg.
\emph{On the prime zeta function.}
BIT Numerical Mathematics 8 (1968), 187--202.
\end{thebibliography}
